\documentclass[11pt,a4paper]{article}

\usepackage[T1]{fontenc}
\usepackage[utf8]{inputenc}
\usepackage[margin=2.4cm]{geometry}
\usepackage{amsmath,amssymb,amsthm}
\usepackage{booktabs}
\usepackage{longtable}
\usepackage{array}
\usepackage{xcolor}
\usepackage{titlesec}
\usepackage{enumitem}
\usepackage[hidelinks,colorlinks=true,linkcolor=BlueLink,urlcolor=BlueLink]{hyperref}
\usepackage{fancyhdr}
\usepackage{tcolorbox}
\tcbuselibrary{skins,breakable}

\definecolor{BlueLink}{HTML}{1A4F8A}
\definecolor{GoodGreen}{HTML}{1B6B3A}
\definecolor{WarnAmber}{HTML}{9A6100}
\definecolor{BadRed}{HTML}{9B2226}
\definecolor{SoftGrey}{HTML}{F4F4F2}
\definecolor{RuleGrey}{HTML}{BBBBBB}

\newcommand{\done}{\textcolor{GoodGreen}{\textbf{DONE}}}
\newcommand{\partialx}{\textcolor{WarnAmber}{\textbf{PARTIAL}}}
\newcommand{\open}{\textcolor{BadRed}{\textbf{OPEN}}}
\usepackage[htt]{hyphenat}
\newcommand{\code}[1]{\texttt{\small\hyphenchar\font=`\-\relax #1}}

\titleformat{\section}{\Large\bfseries}{\thesection}{0.7em}{}[\vspace{-0.6em}\rule{\textwidth}{0.8pt}]
\titleformat{\subsection}{\large\bfseries}{\thesubsection}{0.7em}{}

\pagestyle{fancy}
\fancyhf{}
\lhead{\small BP for diffusion scores --- research compendium}
\rhead{\small\thepage}
\renewcommand{\headrulewidth}{0.4pt}

\newtcolorbox{keybox}[1][]{
  breakable, colback=SoftGrey, colframe=RuleGrey, boxrule=0.5pt,
  left=6pt, right=6pt, top=5pt, bottom=5pt, #1}

\setlist[itemize]{leftmargin=1.4em, itemsep=2pt, topsep=3pt}
\setlist[enumerate]{leftmargin=1.6em, itemsep=2pt, topsep=3pt}

\title{\vspace{-1.5cm}\textbf{Belief propagation for diffusion scores}\\[0.3em]
\large A complete record of the research: questions, answers, code, and what is missing}
\author{Compiled for Giovanni Mantovani\\[0.2em]
\small Advisors: Marc M\'ezard (supervisor), J\'er\^ome Garnier-Brun}
\date{\today}

\begin{document}
\maketitle
\thispagestyle{fancy}

\begin{keybox}
\textbf{Purpose.} This document exists because the research spread across several
packages, two machines and a branch that was never merged, and the thread became hard to
hold. It collects \emph{every} question posed by J\'er\^ome in the call of 29 July 2026 and
by Marc in the follow-up email, together with every question the project posed to itself,
and states for each one what was done analytically, what was measured, where the code and
outputs live, and what is still missing.

\medskip
\textbf{Reading order if you have twenty minutes.} \S\ref{sec:orient} (where everything
lives), \S\ref{sec:arc} (the story in one page), \S\ref{sec:scoreboard} (the scoreboard),
\S\ref{sec:missing} (what is missing).
\end{keybox}

\tableofcontents
\newpage

%======================================================================
\section{Orientation: where everything actually lives}
\label{sec:orient}

This section exists first because it is the part that was lost.

\subsection{The repository and the branch}

\begin{center}
\begin{tabular}{@{}ll@{}}
\toprule
Remote & \url{https://github.com/GioviManto/diffusion} \\
Local clone & \code{/Users/gloriabagnato/Code/Thesis/Diffusion} \\
\textbf{Research branch} & \code{claude/em-bp-denoiser-learning-e07ike} \\
Default branch & \code{main} --- 49 commits \emph{behind} the research branch \\
\bottomrule
\end{tabular}
\end{center}

\begin{keybox}
\textbf{The single most important navigational fact.} Essentially all of the research
lives on \code{claude/em-bp-denoiser-learning-e07ike} and has \textbf{never been merged
into \code{main}}. The branch is 49 commits ahead; \code{main} contains nothing the branch
lacks. If you check out \code{main} --- which is the default --- you will not see the EM
work, the neural baselines, the reverse-diffusion code, the hierarchical experiments, or
any of the reports.

\medskip
To see the work:
\code{git fetch origin \&\& git checkout claude/em-bp-denoiser-learning-e07ike}
\end{keybox}

\subsection{Map of the branch}

\begin{center}
\begin{tabular}{@{}p{6.4cm}p{8.8cm}@{}}
\toprule
\textbf{Path} & \textbf{What is there} \\
\midrule
\code{research/nongaussian-bp/} & \textbf{The main package.} Layers 2--6: all 15
experiments, the EM implementation, the neural baselines, the reverse SDE, the hierarchy
code, 101 tests. \\
\code{\quad src/} & 20 modules --- \code{em.py}, \code{kernels.py}, \code{denoiser.py},
\code{reverse.py}, \code{nnet.py}, \code{hierarchy.py}, \code{bp\_grid.py},
\code{bp\_mixture.py}, \code{discrete.py}, \code{local\_head.py}, \dots \\
\code{\quad experiments/} & \code{exp\_01} \dots \code{exp\_15}, each self-contained and
documented in its own docstring. \\
\code{\quad outputs/} & One directory per experiment; CSVs and PNGs, all committed. \\
\code{\quad tests/} & 10 test files, 101 tests. \\
\code{\quad report/} & \code{em\_bp\_learning.pdf} (14\,pp, five propositions);
\code{layer6\_session\_report.pdf} (15\,pp --- the Layer-6 account, including its own
corrections section); \code{updated\_report.pdf}. \\
\code{\quad hpc/} & SLURM templates --- \textbf{never adapted, never run}. See \S\ref{sec:hpc}. \\
\code{\quad tools/} & \code{merge\_replicates.py}; \code{check\_reproducible.py}, written
after the \code{rng\_for} bug (\S\ref{sec:corrections}) to verify that a fresh interpreter
reproduces a run bit-for-bit. \\
\midrule
\code{questions/}\newline\code{\quad QUESTIONS\_AND\_ANSWERS.md} & \textbf{798 lines.} Every
question from every source, with status. The primary source for \S\ref{sec:questions} of
this document. \\
\code{docs/}\newline\code{\quad EM\_BP\_LEARNING\_COMPENDIUM.md} & 872 lines. The Layer-5/6
results in detail. \\
\code{docs/RESULT\_LEDGER.md} & Every thesis claim classified
EXACT / NUMERICAL / HEURISTIC / SPECULATIVE / FUTURE. \\
\code{docs/PAPER\_CONNECTIONS.md} & Relation to the advisor papers: how they were read, two
corrections that followed, the derivations in this package's conventions, and a ranked
untried list (reproduced in \S\ref{sec:missing}). \\
\code{docs/AGENT\_HANDOFF\_EM\_BP.md} & Practical notes for resuming work. \\
\midrule
\code{research/gaussian-ar1-bp/} & Layer 1--2, the Gaussian AR(1) baseline and the report
that was discussed in the July call. \\
\code{research/bp-from-scratch/}, \code{research/gaussian-bp/},
\code{research/unified-note/} & Earlier analytic packages; \code{numerical\_audit.py} in
the latter two (72/72 and 55/55 passing). \\
\code{thesis/} & The thesis itself --- \code{main.tex}, 9 chapters, figures, bibliography. \\
\bottomrule
\end{tabular}
\end{center}

\subsection{The HPC --- and an important distinction}
\label{sec:hpc}

\begin{keybox}
\textbf{The branch's experiments all ran locally.} \code{research/nongaussian-bp/hpc/}
contains four SLURM array scripts (\code{slurm\_exp06}, \code{slurm\_exp07},
\code{slurm\_layer6}, \code{slurm\_replicates} --- the last a 32-task array) plus a README.
Its own README states they are \textbf{templates}: the lines marked \code{\# SITE} need the
cluster's account, partition and module names, and \textbf{those were never filled in}.
There is no \code{nongaussian-bp} directory on the Bocconi cluster.

\medskip
So every number in \S\ref{sec:questions} --- the 6-seed headline, the 12-replicate rate
study, the 4-seed architecture comparison, all of Layer 6 --- was computed \textbf{on the
laptop}. That is worth knowing both because it bounds what was affordable, and because the
cluster path is available and unused if any of it needs scaling up.
\end{keybox}

\noindent
Separately, there \emph{is} a live cluster project --- but it belongs to the
\code{research/em-learning/} re-implementation (\S\ref{sec:duplication}), not to this
branch:

\begin{center}
\begin{tabular}{@{}ll@{}}
\toprule
Host & \code{lnode01-da.hpc.unibocconi.it} \\
User & \code{3164542} \\
Scheduler & SLURM. QOS caps: \textbf{30 submitted, 10 running} per user. \\
Disk & 144\,G used of a 180\,G soft / 200\,G hard quota. \\
\midrule
\code{\textasciitilde/bp-em-diffusion/} & 175\,M. Dedicated project directory. \\
\code{\quad .venv/} & Dedicated venv: Python 3.13.11, numpy 2.3.5, matplotlib 3.10.7. \\
\code{\quad code/} & The sweep code (a \emph{separate}, smaller re-implementation --- see
\S\ref{sec:duplication}). \\
\code{\quad outputs/} & Three completed sweeps, merged. \\
\code{\quad logs/}, \code{STATUS.md} & Job logs and a status note. \\
\bottomrule
\end{tabular}
\end{center}

\noindent
The three cluster sweeps are \textbf{complete}: 5{,}530 rows, \textbf{zero failed cells},
48 CPU-hours. Merge and retrieve with:

\begin{keybox}
\code{ssh lnode01-da.hpc.unibocconi.it}\\
\code{\textasciitilde/bp-em-diffusion/.venv/bin/python \textasciitilde/bp-em-diffusion/code/merge\_shards.py \textbackslash}\\
\code{\quad --root \textasciitilde/bp-em-diffusion/outputs}\\[0.4em]
\code{rsync -az 'lnode01-da.hpc.unibocconi.it:bp-em-diffusion/outputs/*/*\_all.csv' outputs/}
\end{keybox}

\noindent
Lines beginning \code{!!} in the merge output mean missing or failed cells. On the last run
every line began \code{OK}.

\subsection{A duplication you should know about}
\label{sec:duplication}

There are \textbf{two} EM implementations in the project, and this is worth understanding
because it is confusing otherwise.

\begin{center}
\begin{tabular}{@{}p{4.6cm}p{5.1cm}p{5.5cm}@{}}
\toprule
& \textbf{\code{research/nongaussian-bp/}} & \textbf{\code{research/em-learning/}} \\
\midrule
Where & the research branch & untracked, local only \\
Status & \textbf{authoritative} & superseded \\
Scope & 15 experiments, 4 kernel rungs, neural baselines, reverse SDE, hierarchies & 3 modules, 1 kernel rung, one MLP \\
Tests & 101 & 4 suites \\
Contributes & everything & the 5{,}530-row cluster sweep \\
\bottomrule
\end{tabular}
\end{center}

\noindent
\code{research/em-learning/} was built from scratch in a later session without noticing the
existing package. It reproduces a subset of the same physics with a different
implementation. Its one durable contribution is the large cluster sweep, whose value is
that it \textbf{independently replicates} the branch's central correction on separate code
(\S\ref{sec:cluster}). Everything else in it should be regarded as superseded.

%======================================================================
\newpage
\section{The story in one page}
\label{sec:arc}

\subsection{The scientific problem}

In generative diffusion, a model is trained to denoise. The optimal denoiser is the
posterior mean $\mathbb{E}[a \mid x, t]$, and the score follows from it by Tweedie's
identity
\[
s_i(x,t) \;=\; -\,\frac{x_i - \alpha_t\, \mathbb{E}[a_i \mid x]}{\Delta_t},
\qquad \alpha_t = e^{-t}, \quad \Delta_t = 1 - e^{-2t}.
\]
The difficulty that motivates everything: \textbf{for essentially no interesting data
distribution is the exact score computable}. The field therefore cannot separate
``the network learned the wrong score'' from ``the score itself is hard here''.

There is a second, sharper motivation which J\'er\^ome supplied in the call and which
belongs at the front of the thesis. The global minimiser of the diffusion training loss is
the \emph{empirical} score --- a Gaussian mixture centred on the training points. A model
that attained it would reproduce its training set exactly. That diffusion models
nevertheless generalise is the central open puzzle, and studying it requires data
distributions where the exact score is known.

\subsection{The project's answer}

Take the clean data to be a \textbf{first-order Markov chain}. Then:

\begin{enumerate}
\item Coordinatewise OU noising leaves the posterior a chain --- the likelihood factorises
      sitewise and adds no edges.
\item A chain is a tree.
\item Sum-product belief propagation is \emph{exact} on a tree.
\end{enumerate}

\noindent
So BP computes the exact diffusion score for a non-trivial, non-Gaussian, correlated data
model. That is the asset the whole project is built on, and it is what J\'er\^ome called
``the most important argument'' --- an efficient exact score where the field has almost none.

\subsection{The six layers}

\begin{center}
\begin{tabular}{@{}p{1.1cm}p{6.0cm}p{8.1cm}@{}}
\toprule
& \textbf{What it established} & \textbf{Where} \\
\midrule
L1 & Gaussian AR(1): exact score in closed form; BP $=$ Kalman + RTS smoother. & \code{gaussian-ar1-bp/}, \code{gaussian-bp/} \\
L2 & Grid BP as a calibrated numerical ground truth; the Gaussian closure characterised. & \code{exp\_01}, \code{exp\_02} \\
L3 & Non-Gaussian innovations: when and how the Gaussian closure fails. & \code{exp\_03}, \code{exp\_09}, \code{exp\_11} \\
L4 & Approximate Markovianity: the non-Markov correction is exactly rank one. & \code{exp\_04} \\
L5 & \textbf{Learning.} The kernel becomes unknown and is estimated by EM through exact BP. & \code{exp\_06}--\code{exp\_08}, \code{exp\_10}, \code{exp\_12} \\
L6 & Hierarchical priors, speciation ladders, memorisation. & \code{exp\_13}--\code{exp\_15} \\
\bottomrule
\end{tabular}
\end{center}

\noindent
Layers 1--4 assume the prior is \emph{known}. Layer 5 is the one Marc asked for: remove that
assumption.

%======================================================================
\newpage
\section{The scoreboard}
\label{sec:scoreboard}

Every question from every source, in one table. Details in \S\ref{sec:questions}.

\begin{center}
\small
\begin{longtable}{@{}p{1.0cm}p{9.6cm}p{4.4cm}@{}}
\toprule
\textbf{ID} & \textbf{Question} & \textbf{Status} \\
\midrule
\endhead
\multicolumn{3}{@{}l}{\textbf{[E] Marc's email, 30 July 2026}}\\
\midrule
E1 & Introduce learning: regress BP parameters by EM instead of training a network & \done \\
E2 & ``This does not require noising at all, I think'' & \done\ (sharpened) \\
E3 & Efficient denoiser learning by EM vs a vanilla network $\Rightarrow$ publishable & \done \\
E4 & Start an Overleaf write-up in paper format & \partialx \\
E5 & Does this require autodiff through BP? & \done\ (no) \\
\midrule
\multicolumn{3}{@{}l}{\textbf{[C] J\'er\^ome's call, 29 July 2026}}\\
\midrule
C-i & Gaussian BP ``doing something strange, not just mean/variance'' & \done\ --- he was right \\
C-ii & Problems in the definition of Gaussian BP & \done\ --- changes an interpretation \\
C-iii & Re-run both the grid BP and the Gaussian case & \done \\
C-iv & Then do the learning Marc asked for & \done \\
C-v & Match BP's Gaussian closure to $\alpha_t\Sigma_0\Sigma_t^{-1}x$ analytically & \done \\
C-vi & Check probability mass at the grid boundary & \done \\
C-vii & Posterior mean is the better metric, not relative score & \done \\
C-viii & Compare \emph{generated data}, not just scores (reverse diffusion) & \partialx \\
C-ix & Train an MLP on the Laplace chain, compare to BP & \done \\
C-x & How deep must a network be to implement BP on a chain? & \partialx\ (see N2) \\
C-xi & Vary $\rho$ to see the effect of correlation & \done \\
C-xii & Thesis framing: exact vs empirical score & \open\ (not written) \\
\midrule
\multicolumn{3}{@{}l}{\textbf{[R] The report's own ``next steps''}}\\
\midrule
R1 & More trials and larger grids & \done \\
R2 & Explore different $\rho$ & \done\ --- \emph{sign is opposite} to the conjecture \\
R3 & Laplace, Student, mixture, bounded-support innovations & \done \\
R4 & Mixture \emph{message} approximations: how many components? & \done \\
R5 & Do neural message approximators need BP structure? & \partialx \\
\midrule
\multicolumn{3}{@{}l}{\textbf{[F] The result ledger's ``explicitly future'' list}}\\
\midrule
F1 & Mixture-of-Gaussians message closure & \done\ ($=$ R4) \\
F2 & Discrete-alphabet chain, exact vector messages & \done \\
F3 & Chain $+$ global latent; hybrid BP $+$ learned residual & \partialx \\
F4 & Reverse-SDE dynamics under exact / closed / truncated scores & \partialx \\
F5 & Non-Gaussian locality laws & \done \\
\midrule
\multicolumn{3}{@{}l}{\textbf{[L] Questions raised during Layer 5}}\\
\midrule
L1 & Does the literal \emph{gradient} route work? & \done\ --- not one-sided \\
L2 & Is the observed EM behaviour trustworthy? & \done\ --- five cross-checks \\
L3 & Discretization artifacts? & \done\ --- two, both Laplace-only \\
L4 & Can an unknown innovation law be recovered from noisy data alone? & \done\ --- yes \\
\midrule
\multicolumn{3}{@{}l}{\textbf{[P] Questions from the two advisor papers (Layer 6)}}\\
\midrule
P1 & Does a hierarchical prior show a \emph{ladder} of speciation times? & \done\ --- six of them \\
P2 & Does our chain sit in that regime? & \done\ --- \textbf{no}; bounds generalisation \\
P3 & Does memorisation/collapse apply to a BP score? & \done\ --- the axis does not exist \\
P4 & Can EM through exact BP run on a tree? & \done\ --- yes, M-step unchanged \\
P5 & Does the network acquire hierarchy levels in order? & \done\ --- no (data-limited) \\
P6 & Does the EM-BP advantage grow with correlation range? & \partialx\ --- confounded \\
\midrule
\multicolumn{3}{@{}l}{\textbf{[N] New questions opened by the answers}}\\
\midrule
N1 & What scalar governs locality? & \done\ --- negentropy, not kurtosis \\
N2 & Does a \emph{structured} architecture close the EM-BP gap? & \done\ --- most of it, but the rest \emph{grows} \\
\bottomrule
\end{longtable}
\end{center}

\begin{keybox}
\textbf{Summary.} Of 33 questions: \textbf{25 answered}, \textbf{7 partial with the gap
stated}, \textbf{1 open} (the thesis framing, which is writing rather than research).
No original question is unaddressed.
\end{keybox}

%======================================================================
\newpage
\section{Theory: what was proved}
\label{sec:theory}

Five propositions, in \code{research/nongaussian-bp/report/em\_bp\_learning.tex}
$\rightarrow$ \code{.pdf} (14\,pp, compiles clean).

\begin{center}
\small
\begin{tabular}{@{}p{0.5cm}p{7.2cm}p{7.3cm}@{}}
\toprule
\# & \textbf{Statement} & \textbf{Why it matters} \\
\midrule
1 & \textbf{The E-step is exact.} Sitewise likelihood factors reweight node potentials
without adding edges, so the posterior is still a chain; BP returns exact single-site
\emph{and pairwise} marginals. &
In a generic latent-variable model the E-step is approximate and the monotonicity
guarantee dies with it. Here nothing is approximated but the message representation. \\
\addlinespace
2 & \textbf{The whole E-step is one matrix.} Everything compresses into an $M\times M$
matrix $\Xi$ --- the continuum analogue of Baum-Welch's expected transition counts. &
The M-step never touches the data again (cost $O(PM^2)$ regardless of $N$); observations at
different noise levels simply \emph{add} into one $\Xi$; changing kernel family changes only
the M-step. \\
\addlinespace
3 & \textbf{Fisher's identity $\Rightarrow$ no autodiff.}
$\nabla L(\theta) = \langle \Xi(\theta), \nabla \log K_\theta\rangle$ is the exact gradient
of the marginal log-likelihood. &
BP is never differentiated through --- it \emph{supplies} the expectation. Everything is
pure numpy. Verified against finite differences to $\sim 10^{-9}$. \\
\addlinespace
4 & \textbf{Noising never destroys identifiability.}
$\varphi_{t,\theta}(\xi) = \varphi_\theta(\alpha_t\xi)\exp(-\Delta_t\|\xi\|^2/2)$; the
Gaussian factor has no zeros and $\alpha_t>0$ for finite $t$. &
Makes Marc's ``does not require noising'' precise: identifiability survives \emph{any}
finite $t$. What degrades is \emph{information}. Injective but ill-conditioned. \\
\addlinespace
5 & \textbf{Denoiser sensitivity is a posterior covariance.}
$\partial m_i/\partial\theta_p = \mathrm{Cov}(a_i, \sum_l \partial_{\theta_p}\log K(a_l|a_{l-1}))$. &
One $N^{-1/2}$-consistent parameter fit gives an $N^{-1/2}$-accurate denoiser
\emph{uniformly in $t$}. The network gets no such statement --- its error must be paid for
separately in every region of $(x,t)$. \\
\bottomrule
\end{tabular}
\end{center}

\begin{keybox}
\textbf{The structural argument in one sentence.} The learned parameters live on the
\emph{clean} chain; the noise level lives in the \emph{likelihood}. So one fit serves every
$t$, and the risk decomposition carries no function-class approximation term --- given a
prior, BP returns the exact Bayes denoiser.
\end{keybox}

\subsection{The four kernel rungs}

All consume the same $\Xi$. In \code{src/kernels.py}.

\begin{enumerate}
\item \textbf{Gaussian AR(1)} $(\rho,q)$ --- closed-form M-step (belief-weighted
Yule--Walker). The correctness harness, since everything is analytically available.
\item \textbf{Laplace AR(1)} $(\rho,b)$ --- closed form \emph{despite} non-differentiability:
profile out $b$, and maximising over $\rho$ becomes weighted least-absolute-deviations,
solved exactly by a weighted median of grid ratios. \emph{This rung has a lattice-bias
pathology --- see \S\ref{sec:corrections}.}
\item \textbf{Mixture innovation} $(\rho, \{\pi,\mu,\sigma^2\}_c)$ --- the honest ``we know
it is Markov but not the model'' estimator. Has never heard of the Laplace density it must
recover. Closed-form ECM via an inner latent.
\item \textbf{Mixture-density network} --- a small MLP emitting a state-dependent mixture.
By Fisher's identity it is evaluated and backpropagated \emph{only at the $M$ grid points},
once per EM iteration, regardless of dataset size.
\end{enumerate}

%======================================================================
\newpage
\section{The questions, answered in full}
\label{sec:questions}

\subsection{[E] Marc's email}

\subsubsection*{E1. EM instead of a neural network --- \done}

This is the whole of Layer 5. Theory in \code{report/em\_bp\_learning.pdf};
implementation in \code{src/em.py}, \code{src/kernels.py}, \code{src/denoiser.py}.

Four structural points, none obvious at the outset: the E-step is \emph{exact}, not
variational; the entire E-step compresses to one $M\times M$ matrix; \textbf{no autodiff is
needed anywhere}; and one fit serves every noise level, because the learned object lives on
$\mathbb{R}\times\mathbb{R}$ and carries no $t$.

\subsubsection*{E2. ``Does not require noising at all'' --- \done, and sharpened}

Marc's remark conflated two things, and separating them is the useful answer.

\begin{itemize}
\item \textbf{Identifiability is never destroyed by noising} (Prop.~4). The model is
identifiable at any finite $t$ exactly when it is identifiable from clean data.
\item \textbf{Information is destroyed, and fast.} Injective does not mean stable --- this is
Gaussian deconvolution and the inverse map is unbounded. Measured per chain: $J[q,q]$ falls
\textbf{142$\times$} from $t=0.05$ to $t=1.6$ while $J[\rho,\rho]$ falls \textbf{26$\times$}.
\end{itemize}

\noindent
So the remark is right --- clean data is the $t\to 0$ limit and the most informative case,
and \code{fit\_clean} does exactly that with no BP required. But the reason is not that
noise breaks identification; it is that noise costs \emph{information}, and it hits the
innovation \emph{shape} about 5$\times$ harder than the correlation.
\textbf{Non-Gaussianity is the first thing the channel destroys} --- which is precisely the
property this project cares about.

\subsubsection*{E3. The publishable claim --- \done}

Six independent training seeds, extended to $N=4096$. Test set and its exact-BP reference
held fixed across seeds. Relative score error averaged over the noise schedule, $\pm 1$
standard error:

\begin{center}
\begin{tabular}{@{}rrrr@{}}
\toprule
$N$ chains & best network & \textbf{EM-BP} & ratio \\
\midrule
32   & $0.6363 \pm 0.0043$ & $\mathbf{0.0651 \pm 0.0031}$ & 9.8$\times$ \\
128  & $0.5117 \pm 0.0025$ & $\mathbf{0.0333 \pm 0.0035}$ & 15.4$\times$ \\
512  & $0.2813 \pm 0.0013$ & $\mathbf{0.0211 \pm 0.0022}$ & 13.3$\times$ \\
1024 & $0.2154 \pm 0.0009$ & $\mathbf{0.0182 \pm 0.0021}$ & 11.8$\times$ \\
2048 & $0.1756 \pm 0.0009$ & $\mathbf{0.0142 \pm 0.0009}$ & 12.4$\times$ \\
4096 & $0.1643 \pm 0.0009$ & $\mathbf{0.0124 \pm 0.0005}$ & 13.2$\times$ \\
\bottomrule
\end{tabular}
\end{center}

\noindent
EM-BP uses \textbf{13 parameters}; the network \textbf{25{,}248}. EM-BP on \textbf{32}
chains beats the network on \textbf{4096} by a factor of 2.5, with non-overlapping error
bars --- so the network needs \textbf{more than 128$\times$ the data}, and that is a lower
bound, since 32 was the smallest budget tried.

\medskip
\textbf{Three controls make this more than a favourable setup:}
\begin{itemize}
\item \emph{Capacity is not the explanation.} 24 configurations from 3.2k to 297k
parameters; the best network reaches 0.208 against EM-BP's 0.022, and error \emph{degrades}
with capacity.
\item \emph{The yardstick is verified.} The reference denoiser everything is scored against
agrees with 6M-sample importance sampling to $1.1\times 10^{-3}$ relative.
\item \emph{The baseline works.} Trained on a Gaussian chain where the exact denoiser is
closed-form and linear, it reaches 0.099 --- so the gap is a sample-efficiency effect, not a
broken competitor.
\end{itemize}

\begin{keybox}
\textbf{Two things stated against the claim.}
\begin{enumerate}
\item BP inference is \textbf{211--320$\times$ slower} per evaluation. Reverse diffusion
calls the denoiser at every integration step, so this is the weakest point of the story.
\item The baseline is a \emph{vanilla} MLP, as the email specified. A locality-respecting
architecture closes most of the gap --- see N2, where the honest margin is
\textbf{4.1$\times$}, not 10$\times$.
\end{enumerate}
\end{keybox}

\subsubsection*{E4. The Overleaf write-up --- \partialx}

\code{report/em\_bp\_learning.tex} was written paper-shaped on purpose: context, numbered
propositions with proofs, detailed calculation delegated to remarks and appendices,
limitations as a first-class section. It compiles clean to 14\,pp.

\textbf{Gap:} it has \textbf{no results section} --- the numbers live in the compendium and
the CSVs --- and \textbf{no Overleaf project exists}. Converting it is mechanical; \S4 of
the compendium is the raw material. \emph{This is the item with the deadline attached.}

\subsubsection*{E5. Does this need autodiff through BP? --- \done: no}

Advice received elsewhere held that EM here ``will require moving the grid BP
forward-backward passes into an automatic differentiation framework''. That is wrong, and
it is the single most useful practical finding of the layer. Fisher's identity gives the
gradient of the \emph{exact} marginal log-likelihood from one BP pass, with nothing
differentiated through the recursion. BP is not a computation graph to backpropagate; it is
the oracle that supplies the expectation. Verified against finite differences to
$\sim 10^{-9}$. The entire implementation is pure numpy.

\subsection{[C] J\'er\^ome's call}

\subsubsection*{C-i, C-ii, C-iii. The Gaussian BP was wrong --- \done, he was right}

J\'er\^ome's concern in the call was that Gaussian-projected BP was ``not updating just the
mean and variance \dots but doing something strange'', and that Figure~2 ought to beat
Figure~1 because it carries no discretization. He was right, and the bug was worse than a
bug --- \textbf{it invented information}.

The old routine evaluated each Gaussian message \emph{on the grid}, pushed it through the
exact grid update, then moment-matched. At weakly informative $t$ the outgoing message is a
near-flat ramp whose maximiser lies outside $[-A,A]$; moment-matching that \emph{truncated}
function drags the mean toward the boundary, and the next update pushes it further ---
positive feedback, ending with the message pinned at the edge. Equivalently: an infinite
variance became $(2A)^2/12 \approx 21$ on a truncated grid, injecting spurious precision
$\sim 3/A^2$.

The replacement does exactly what the call asked: \textbf{analytic information form},
carrying only precision $\lambda$ and information $h$, closed-form updates, no grid
anywhere. A flat message is $\lambda=0$ exactly, so nothing is ever truncated.

\begin{center}
\begin{tabular}{@{}rrr@{}}
\toprule
$t$ & information form & legacy grid projection \\
\midrule
0.2 & $1.2\times10^{-15}$ & $3.9\times10^{-6}$ \\
0.6 & $1.0\times10^{-15}$ & $9.4\times10^{-3}$ \\
1.0 & $4.4\times10^{-16}$ & $1.1\times10^{-1}$ \\
1.3 & $6.7\times10^{-16}$ & $4.9\times10^{-1}$ \\
1.8 & $4.4\times10^{-16}$ & $\mathbf{9.1\times10^{-1}}$ \\
\bottomrule
\end{tabular}
\end{center}

\noindent
(max absolute error in the posterior mean). Exact to machine precision at every noise level;
the legacy form up to six orders of magnitude worse, and worsening precisely as the
likelihood stops being informative --- the signature of the mechanism above.

\medskip
\textbf{C-ii changes an interpretation.} For \emph{linear-transition} chains with Gaussian
OU likelihoods, moment-matched single-Gaussian BP is \emph{mathematically identical} to
exact Gaussian BP on the covariance-matched AR(1) model: the transition step maps
$N(m,v)$ to something with first two moments $(\rho m, \rho^2 v + q)$ \emph{regardless of
innovation shape}. So ``Gaussian \textbf{message} approximation error'' and ``Gaussian
\textbf{model} approximation error'' are the same object at the single-Gaussian level. What
the old package reported as a message-representation error was a model error. The values
were right; the label was not. This is exactly why R4 (mixture-message closure) was the
open question worth doing --- it is the first place the two are separable.

\textbf{C-iii: both re-run}, under corrected code \emph{and} fixed seeding. Grid BP is
spectrally accurate at the working default: worst relative error $9.2\times10^{-15}$ at
$M=401$, $A=8$. Gaussian BP is clean, monotone, artifact-free, with the score/mean identity
holding to $\le 2.9\times10^{-12}$ on every row.

\subsubsection*{C-v. Match the Gaussian closure to $\alpha_t\Sigma_0\Sigma_t^{-1}x$ --- \done, exceeded}

J\'er\^ome asked whether one could show explicitly that propagating BP's Gaussian messages
to their fixed point reproduces the closed-form Gaussian score. It does, and the result is
stronger than asked: the closure uses \emph{only} the innovation's first two moments, so it
returns the LMMSE denoiser for \textbf{any} additive innovation with mean 0 and variance
$q$, not just Gaussian. Verified to $\sim 6\times10^{-16}$ across five families, and
independently confirmed by the cluster sweep, where the Gaussian column is machine epsilon
(\S\ref{sec:cluster}).

\subsubsection*{C-vi. Boundary mass --- \done}

J\'er\^ome: ``you could add something in the code that checks that the support at the borders
of the grid is small enough.'' Implemented as \code{boundary\_mass()} and recorded per cell.
Grid widened $[-8,8]\to[-12,12]$; boundary mass $7.4\times10^{-7} \to 5.3\times10^{-10}$.
Worst value over all 160 sweep cells: $2.0\times10^{-5}$. Truncation drives nothing.

\subsubsection*{C-vii. Posterior mean is the better metric --- \done}

Adopted throughout. The denoising loss is additionally framed against its \emph{Bayes-risk
floor}, since raw loss reduction is meaningless --- an earlier test demanded a $>40\%$
reduction when the arithmetic maximum was 39\%.

\subsubsection*{C-viii. Compare generated data, not scores --- \partialx}

This is the point J\'er\^ome returned to most often: ``the error on the score \dots and the
distribution of data you end up generating \dots is not trivial, because you accumulate the
errors of the scores, and not all times are equal.''

\code{exp\_05} ran it. Under the Gaussian score on a Laplace chain: second-order statistics
are reproduced, but heavy tails are \textbf{washed out} --- excess kurtosis \textbf{0.12}
against a true \textbf{2.7--2.9}. The score error is dynamically stable in $L^2$ (trajectory
divergence 0.11 despite pointwise deviation 0.49 at small $t$) yet \emph{distributionally
decisive} in higher moments. That is the sharpest available statement of why score error and
sampling error are different quantities.

\textbf{Gap:} never run with a \emph{learned} score. Now that Layer 5 produces one, the
obvious experiment is reverse dynamics under the EM-learned kernel versus the network ---
which is also where BP's 211--320$\times$ inference penalty would actually bite.

\subsubsection*{C-ix, C-x. MLP vs BP, and the depth question --- \done\ / \partialx}

J\'er\^ome asked for the Laplace chain specifically, and for good reason: in the Gaussian
case BP is linear and he called Markovianity there ``completely useless \dots basically like
doing covariance''. \code{exp\_07} does exactly this on the Laplace chain (E3 above).

On depth --- ``the biggest question at the moment is what is the required depth of a neural
network that you need to learn it?'' --- \code{exp\_12} answers the adjacent and more
tractable question of \emph{receptive field} rather than depth. See N2 and
\S\ref{sec:receptive}. The literal depth question, in the form J\'er\^ome posed it (can a
narrow-but-deep network implement the two BP sweeps, and how many layers does it need?),
\textbf{has not been tested}.

\subsubsection*{C-xi. Vary $\rho$ --- \done, and the sign is a surprise}

See R2 below. Stronger correlation \emph{helps} the Gaussian closure, contrary to the
conjecture in the report.

\subsubsection*{C-xii. Thesis framing --- \open}

J\'er\^ome handed over the whole narrative arc in the call: the empirical-score mystery
motivates needing exact scores on controllable data $\to$ Markov chains are such a model
$\to$ BP computes the exact score $\to$ compare against trained networks. \textbf{Not a line
of it is written down.} This is writing, not research, and it is on the critical path.

\subsection{[R] The report's own next steps}

\subsubsection*{R1. More trials and larger grids --- \done}

Grid: relative error against an $M=1201$ reference is $1.9\times10^{-5}$ ($M=201$),
$4.1\times10^{-6}$ ($M=401$), $1.5\times10^{-6}$ ($M=601$). $M=401$, $A=8$ is the working
default, four orders of magnitude below any learning error in the project.

Trials: this mattered more than expected. The $N^{-1/2}$ rate study at \textbf{4} replicates
gave slopes $-0.26, -0.69, -0.25, -0.41$ against a predicted $-0.50$ and did \emph{not}
support the rate. Re-run at \textbf{12} replicates it gives a combined slope of
$\mathbf{-0.500 \pm 0.048}$, excluding $-0.25$ and $-0.75$ at $5.2\sigma$.

\begin{keybox}
\textbf{Rule adopted:} any rate claim needs replicates in the \emph{tens}, not the units.
An RMSE over 4 samples cannot resolve a slope better than about $\pm 0.2$.
\end{keybox}

\subsubsection*{R2. Explore different $\rho$ --- \done, sign opposite to the conjecture}

\begin{center}
\begin{tabular}{@{}lrrr@{}}
\toprule
family & $\rho=0.5$ & $\rho=0.85$ & $\rho=0.95$ \\
\midrule
Laplace          & 0.520 & 0.354 & 0.270 \\
Student-$t$(5)   & 0.317 & 0.241 & 0.135 \\
mixture $\kappa=0.9$ & 0.939 & 0.714 & 0.368 \\
\bottomrule
\end{tabular}
\end{center}

\noindent
(median relative score error at small $t$.) The report conjectured stronger correlation
would make messages \emph{more} non-Gaussian and so harder. The measurement says the
reverse: \textbf{a more informative message is closer to Gaussian, because it concentrates}.

\subsubsection*{R3. Four innovation families --- \done, with a genuine negative}

Laplace, Student-$t$ ($\nu=5,8$), symmetric Gaussian mixture ($\kappa=0.3,0.6,0.9$),
spanning excess kurtosis $-1.62$ to $+6.00$; plus bounded support (\code{UniformAR1},
excess kurtosis $-1.20$).

The bounded case was the interesting one because it fails a Gaussian closure through
\emph{support} rather than shape --- the true posterior is exactly zero outside a finite
interval and no Gaussian message can represent that edge. The question was whether
$|$excess kurtosis$|$ survives as the one-number summary when the failure mode changes in
kind. \textbf{It does.} Fitting
$\log_{10}(\text{error}) = 0.801|\text{kurt}| - 2.234$ on the three mixtures alone and
extrapolating to the uniform prior predicts $5.33\times10^{-2}$; measured
$\mathbf{6.07\times10^{-2}}$, a factor of 1.14. A qualitatively different failure mode did
\emph{not} produce a qualitatively different error.

\subsubsection*{R4 / F1. Mixture-message closure --- \done}

Messages are Gaussian mixtures; multiply, forward-push and backward-push are exact and
closed; component count multiplies at every step, so Runnalls' KL-based pairwise merge
collapses back to $C$. \textbf{That merge is the only approximation in the recursion.}

\begin{center}
\begin{tabular}{@{}lrr@{}}
\toprule
prior & $C$ for 10$\times$ over single-Gaussian & $C$ for 100$\times$ \\
\midrule
$\kappa=0.3$ (kurt $-0.18$) & 2 & 4 \\
$\kappa=0.6$ (kurt $-0.72$) & 2 & 4 \\
$\kappa=0.9$ (kurt $-1.62$) & 3 & 6 \\
\bottomrule
\end{tabular}
\end{center}

\noindent
\textbf{Two to three components buy an order of magnitude; four to twelve buy two.}

And --- the separation the project wanted since C-ii --- on a Laplace chain, which the
mixture family \emph{cannot} represent exactly, the curve flattens onto a floor. That floor
is \textbf{model} error; everything above it is \textbf{representation} error:

\begin{center}
\begin{tabular}{@{}lrr@{}}
\toprule
kernel fitted with & $C=1$ & floor reached \\
\midrule
2 components & $5.28\times10^{-2}$ & $9.25\times10^{-3}$ \\
4 components & $5.44\times10^{-2}$ & $3.33\times10^{-3}$ \\
8 components & $5.76\times10^{-2}$ & $1.68\times10^{-3}$ \\
\bottomrule
\end{tabular}
\end{center}

\noindent
This is the first place in the project where the two error sources are visible as separate
quantities rather than the single conflated number.

\subsubsection*{R5. Do neural message approximators need BP structure? --- \partialx}

Three data points, pointing the same way but with a twist:

\begin{itemize}
\item \emph{Layer 4:} a residual MLP on a \emph{frozen} Markov-BP score beats a direct score
MLP by ${\sim}20\times$. Structure preserved $\to$ large gain.
\item \emph{Layer 5:} EM-learned BP beats a vanilla score network by ${\sim}10\times$.
Structure preserved $\to$ large gain.
\item \emph{The other direction:} a mixture-density network placed \emph{inside} the BP
recursion --- maximal structure preservation --- is \textbf{dominated} by a 13-parameter
mixture kernel (held-out score error 0.052 vs 0.026).
\end{itemize}

\noindent
So the sharper reading is that \textbf{structure preservation is necessary, not
sufficient}: rung 4 preserves it perfectly and still loses, because the extra flexibility
buys nothing and costs optimisation difficulty. \emph{Gap:} all three use chain priors where
a low-dimensional family is adequate.

\subsection{[F] The ledger's future list}

\subsubsection*{F2. Discrete alphabet --- \done, and it is the cleanest result in the project}

\code{src/discrete.py}, \code{exp\_10}. The clean chain takes finitely many values; noising
is still the continuous OU channel, so observations are real-valued, but messages are
$S$-vectors. \textbf{BP is exact up to roundoff} --- verified against explicit enumeration
of all $S^n$ configurations: $\Xi$ to $4.4\times10^{-15}$, posterior means to
$2.9\times10^{-15}$, log-evidence to \emph{exactly} zero difference.

Every discretization caveat disappears simultaneously: no quadrature, no truncation, no
resolution condition, and in particular no ratio-lattice quantization.

\textbf{The Baum-Welch analogy becomes an identity.} Here $\Xi$ \emph{is} the expected
transition-count matrix and the M-step is exactly Baum-Welch's: normalise the counts.
Monotonicity violation is \textbf{exactly 0.0}.

\textbf{A confound-free retest of the headline}, with the reference exact rather than a
fine-grid approximation ($S=5$, 20 learned parameters vs 25{,}248):

\begin{center}
\begin{tabular}{@{}rrrr@{}}
\toprule
$N$ & best network & \textbf{EM-BP} & ratio \\
\midrule
32   & 0.761 & \textbf{0.184} & 4.1$\times$ \\
128  & 0.641 & \textbf{0.085} & 7.5$\times$ \\
512  & 0.398 & \textbf{0.052} & 7.7$\times$ \\
2048 & 0.263 & \textbf{0.032} & 8.3$\times$ \\
\bottomrule
\end{tabular}
\end{center}

\noindent
EM-BP on 32 chains again beats the network on 2048, reproducing the $\ge 64\times$ gap with
\textbf{zero grid error}.

\begin{keybox}
\textbf{Two honest findings that qualify the headline.}

\emph{The advantage erodes as the model grows.} At $N=512$, the ratio over the best network
falls $15.4\times$ ($S=3$, 6 params) $\to 17.7\times$ ($S=4$) $\to 8.0\times$ ($S=6$)
$\to 5.7\times$ ($S=8$) $\to 4.4\times$ ($S=12$, 132 params). Still ahead, but the trend is
unambiguous: \textbf{the sample-efficiency advantage is tied to the parametric dimension
being small}, exactly as Prop.~5 predicts. Nothing here says a structured estimator wins
once its parameter count approaches the network's.

\emph{The recovery rate degrades with $S$, and it is not a metric artifact.} Slopes go
$-0.55$ ($S=3$) to $-0.18$ ($S=8$); re-run with a per-entry metric the degradation persists.
The mechanism is \textbf{level crowding}: at fixed unit variance the $S$ levels sit
${\sim}1/S$ apart while the channel resolves only $\sqrt{\Delta_t}/\alpha_t$. At $S=8$ that
ratio is already $0.93 < 1$ at the \emph{lowest} training noise --- adjacent levels are not
distinguishable through the channel at all. An identifiability limit imposed by the channel,
not a failure of the estimator.
\end{keybox}

\subsubsection*{F3. Chain $+$ global latent --- \partialx}

Layer 4 established the operator-level result: for AR(1)-plus-global-latent priors the
non-Markov score correction is \textbf{exactly rank one} (Woodbury), and a residual MLP on a
frozen Markov-BP score is ${\sim}20\times$ more efficient than a direct score net.

\textbf{Gap:} that used the \emph{true} chain prior. A \emph{learned} chain kernel carrying
a learned low-rank correction is untouched. This is the single most natural next piece of
research the repository contains.

\subsubsection*{F5. Non-Gaussian locality --- \done}

Six families at matched $(\rho,q)$, so they differ only beyond second moments. Ratio of
fitted locality decay rate to the Gaussian rate:

\begin{center}
\begin{tabular}{@{}lrrrrr@{}}
\toprule
$\rho$ & $t=0.05$ & $t=0.10$ & $t=0.20$ & $t=0.40$ & $t=0.80$ \\
\midrule
0.50 & \textbf{2.08} & 1.65 & 1.19 & 1.05 & 0.99 \\
0.85 & \textbf{1.47} & 1.19 & 1.06 & 1.02 & 1.00 \\
0.95 & 1.10 & 1.01 & 1.01 & 1.00 & 1.00 \\
\bottomrule
\end{tabular}
\end{center}

\noindent
Non-Gaussian chains are \textbf{less local} than their covariance-matched Gaussian, and the
excess vanishes in both limits --- largest at low noise and weak correlation, gone by
$t\gtrsim 0.4$ or $\rho\gtrsim 0.95$.

\textbf{Architectural consequence:} since $r \sim \log(1/\epsilon)/\log(1/q)$, a receptive
field sized on Gaussian intuition is too small by up to $\mathbf{2.45\times}$ for bimodal
data at low noise.

\subsection{[L] Questions raised during Layer 5}

\subsubsection*{L1. Does the literal gradient route work? --- \done, not one-sided}

Marc's email described \emph{gradient} ascent, not EM, so both were measured. They share the
E-step and $\Xi$ and differ only in what they do with it. \textbf{The split is not the one
intuition suggests:}

\begin{itemize}
\item \textbf{Smooth (Gaussian) kernel:} gradient ascent converges to \emph{exactly} EM's
optimum (log-likelihoods matching to $2\times10^{-9}$ nats at $\eta=0.5$). At $\eta=2$ it
diverges (527 nats of monotonicity violation); at $\eta=8$ it leaves the admissible set.
\textbf{Nothing gained over EM, and a step size must be tuned to lose nothing.}
\item \textbf{Laplace kernel:} gradient ascent is genuinely \textbf{better} --- higher
likelihood ($-8064.7$ vs $-8066.6$) and better on both parameters. Because there EM's M-step
is exactly optimal \emph{for the discretized model} and lands on a lattice point (L3).
\end{itemize}

\noindent
\textbf{Recommendation:} exact M-step where it exists \emph{and} the kernel is smooth;
gradient route where the M-step has no closed form or where its exact solution is a
discretization artifact.

\subsubsection*{L2. Is the EM behaviour trustworthy? --- \done, five independent cross-checks}

These deliberately take routes sharing no code with what they test.

\begin{center}
\small
\begin{tabular}{@{}p{5.4cm}p{6.0cm}p{3.4cm}@{}}
\toprule
\textbf{check} & \textbf{route} & \textbf{result} \\
\midrule
$\Xi$ vs brute force & enumerate all $M^n$ discretized configurations and marginalise by
summation --- no messages anywhere & $\mathbf{1.0\times10^{-14}}$ \\
Evidence vs Monte Carlo & importance sampling on the \emph{Laplace} chain, where no closed
form exists & $\mathbf{0.59}$ s.e. (4M samples) \\
EM fixed point is stationary & Fisher's identity, independent of how the M-step chose its
update & $\|\nabla L\|$ falls $\mathbf{8.1\times10^{6}}\times$ \\
Reference denoiser vs importance sampling & the yardstick every exp\_07 number is measured
against & $\mathbf{1.1\times10^{-3}}$ relative \\
Baseline network can learn & trained on a Gaussian chain with a closed-form linear target &
best error $\mathbf{0.099}$ \\
\bottomrule
\end{tabular}
\end{center}

\noindent
Plus \textbf{monotonicity violation of exactly 0} in every EM run across every experiment ---
the sharpest single test, since an error anywhere in the recursion, the accumulation or the
M-step generically breaks it.

\subsubsection*{L3. Discretization artifacts --- \done: two, both Laplace-only}

\textbf{Lattice quantization.} The Laplace M-step minimises $\sum \Xi|u_k - \rho u_j|$, whose
minimiser is a \emph{breakpoint} --- one of the ratios $u_k/u_j$, i.e.\ a rational. On a
uniform grid through the origin, low-denominator values like $4/5$ have many aliases
($8/10$, $12/15$, \dots) pooling weight onto them. \textbf{$4/5$ is an attractor.} For
$\rho^*=0.7913$, chosen deliberately off the lattice, $\hat\rho$ is pinned at exactly
$0.8000$ across an \textbf{8$\times$ grid refinement}, holding a constant $0.0087$ bias. It
is a \emph{bias}, not a resolution limit --- more grid points do not help. It contaminates
$b$ by 3--4$\times$ where $\rho$ is snapped.

\textbf{Gradient quadrature.} $\partial_\rho \log K$ carries $\mathrm{sign}(e)$; trapezoidal
quadrature of a discontinuous integrand loses the spectral accuracy the rest of the package
enjoys --- off ${\sim}13\%$ at $M=201$, still ${\sim}0.1\%$ at $M=1601$. The $b$ direction,
being smooth, is exact to machine precision --- confirming Prop.~3 is fine and the
quadrature is not.

\begin{keybox}
\textbf{Consequence adopted:} every rate and accuracy claim in the project uses the Gaussian
or mixture kernel, \textbf{never the Laplace one}.
\end{keybox}

\subsubsection*{L4. Can an unknown innovation law be recovered from noisy data? --- \done: yes}

The mixture kernel has never heard of the Laplace density it must recover. Fitted to
\textbf{noisy observations only} ($N=1024$, one noisy realisation per chain, noise up to
$t=1.6$, never a clean sample):

\begin{center}
\begin{tabular}{@{}lrrrr@{}}
\toprule
$C$ & $\hat\rho$ & innov.\ var & excess kurtosis & $\log L$ \\
\midrule
2 & 0.806 & 0.363 & 1.95 & $-42659.4$ \\
3 & 0.807 & 0.360 & 2.23 & $-42658.4$ \\
5 & 0.808 & 0.360 & 2.62 & $-42655.7$ \\
8 & 0.807 & 0.361 & \textbf{2.71} & $-42655.5$ \\
\emph{true} & \emph{0.800} & \emph{0.360} & \emph{3.00} & \\
\bottomrule
\end{tabular}
\end{center}

\noindent
Variance recovered essentially exactly from $C=3$; the heavy tail climbs monotonically
toward its true value; monotonicity violation exactly 0 throughout.

\textbf{Caveat, recorded because it was first got wrong:} the \emph{kurtosis} estimate is
high-variance. At one replicate per $N$ it reads $2.91, 3.45, 3.70, 3.29, 2.64$ for
$N=128\dots2048$ --- non-monotone, and overshooting the true 3.0. What improves monotonically
is the quantity that matters, the denoiser error ($0.083 \to 0.037$).

\subsection{[P] Questions from the two advisor papers}

\begin{keybox}
\textbf{Access status, stated precisely --- this changed partway through the work.}

\medskip
\textbf{P1 (hierarchical filtering, arXiv:2408.15138) has been read in full.} Direct HTTP is
blocked at the environment's gateway (403 on arXiv, Nature, PMC, HAL, PMLR), but the Hugging
Face Hub exposes \code{hf://papers/2408.15138/paper.md} through a server-side route
returning the complete LaTeXML rendering including appendices. Everything about P1 is from
the source.

\medskip
\textbf{P2 (dynamical regimes, Nat.\ Commun.) has \emph{not} been read.} It is not on that
route and every direct route is blocked. What is said about it comes from web search at
abstract level, and \textbf{nothing mathematical is quoted from it}. The speciation crossover
used throughout is derived from scratch in this package's conventions and checked against
40{,}000 sampled forward trajectories. \emph{One search summary rendered the condition with a
reciprocal, $t_S = \tfrac12\log(1/\Lambda)$, contradicting its own statement that $t_S$
diverges with dimension --- deriving it independently avoided importing that slip.}
\textbf{Read P2 before citing it.}
\end{keybox}

\noindent
\textbf{Reading P1 overturned two substantive things}, both corrected in place:

\begin{enumerate}
\item \textbf{The filtering construction was wrong.} ``Filter at level $k$'' had been
implemented as $b^k$ \emph{independent} subtrees. P1 §2.2 instead draws the depth-$k$ nodes
\textbf{conditionally independently given the root}, with the marginals the unfiltered model
gives them. The blocks therefore stay correlated \emph{through the root}; at $k=L$ the leaves
become conditionally i.i.d.\ given the root --- the regime where P1 notes a Naive Bayes
classifier is optimal and attention is superfluous. Independent blocks would have put a zero
where the correct construction puts $\rho^{2L}$. A test now pins the distinction.

\item \textbf{The claim that the hierarchy\,$\times$\,diffusion intersection was
``untouched by either paper'' was false.} Sclocchi, Favero and Wyart, \emph{A phase
transition in diffusion models reveals the hierarchical nature of data}, PNAS 122(1)
e2408799121 (2025), arXiv:2402.16991 --- \textbf{cited by P1 itself} --- studies exactly
diffusion on hierarchically generated data. The measurements here turn out to be
\emph{consistent with} that picture rather than novel against it.
\end{enumerate}

\begin{keybox}
\textbf{Consequence for the thesis:} that third paper is the closest prior work and must be
read before anything about hierarchy and diffusion is written. An earlier version of the
project's own notes claimed its territory was empty.
\end{keybox}

\medskip
\noindent
\textbf{A methodological correction that follows.} P1's probe for sequential learning is
\emph{not} ``resolve the error onto levels and watch it move'' --- which is what
\code{exp\_13 ordering} did, finding nothing. It is to compare the model's predictions
against the \emph{family} of mismatched oracles $\mathrm{BP}_k$ and ask which one it
currently agrees with. \code{exp\_15} implements that; see P5.

\subsubsection*{P1. A ladder of speciation times --- \done: six of them}

The balanced-tree prior has an ultrametric leaf covariance with exactly $L+1$ distinct
eigenvalues, so it has $L+1$ speciation times. On a depth-5 binary tree, $\rho=0.9$:
predicted $t_S = \tfrac12\log(1+\Lambda)$ against the measured crossing of the commitment
curve through $1/\sqrt2$ gives
$1.363/1.368$, $0.707/0.732$, $0.516/0.521$, $0.346/0.340$, $0.202/0.208$, $0.087/0.087$ ---
agreeing to $\le 3.5\%$ at every level. The reverse SDE driven by exact tree BP reproduces
the same ladder and therefore \textbf{resolves the hierarchy coarse-to-fine}.

\emph{Reported rather than trimmed:} the reverse crossing at the finest level lands at 0.125
against 0.087 predicted --- the Euler--Maruyama step and the $t_{\min}=0.02$ floor, which
bite hardest exactly there.

\subsubsection*{P2. Does our chain sit in that regime? --- \done: no, and this bounds the project}

The AR(1) top eigenvalue is bounded by $(1+\rho)/(1-\rho)$ at any length: 64$\times$ the
chain length moves it under 2$\times$ ($\rho=0.85$: $5.49 \to 12.32$ against a limit of
12.33). So its speciation time \textbf{saturates}: a stationary Markov chain has no diverging
speciation time and \textbf{no cascade}. Only the hierarchy produces them.

\begin{keybox}
\textbf{This is the most important limiting result in the project.} Any claim transferred
from this chain model to image-like or hierarchically structured data has to cross that gap
\emph{explicitly}. It should appear in the thesis as a scope statement, not be discovered by
an examiner.
\end{keybox}

\subsubsection*{P3. Does memorisation/collapse apply to a BP score? --- \done: the axis does not exist}

Not a matter of degree. The empirical score's sufficient statistic \emph{is} the training set
--- $N\times n$ numbers --- which is why it must collapse onto that set. A BP score's
sufficient statistic is the fitted kernel, a fixed-size object independent of $N$. There is
no mechanism by which it can memorise, because it never retains the samples. This supplies
the \emph{mechanism} behind the Layer-5 headline: the sample-efficiency advantage and the
absence of memorisation are the same fact seen from two sides.

\subsubsection*{P4. Can EM through exact BP run on a tree? --- \done: yes, M-step unchanged}

\code{tree\_e\_step} produces the \textbf{same $\Xi$} as the chain E-step. That is where the
topology stops being visible, so every kernel consumes it without modification --- a
chain-trained and a tree-trained M-step are the same code. Verified against the closed-form
Gaussian marginal likelihood to relative $10^{-6}$.

\textbf{But convergence is much slower, and this matters practically.}
$\hat\rho = 0.7360 \to 0.7483 \to 0.7487$ at 50/100/150 iterations (true 0.75), zero
monotone violation throughout. Internal nodes are never observed, so the missing-information
fraction is far larger than on a chain. \emph{An estimate read at 40 iterations looks like a
broken M-step and is merely unconverged --- this cost real time.}

\subsubsection*{P5. Does the network acquire hierarchy levels in order? --- \done: no}

\textbf{Two probes, and the second is the right one.} The first
(\code{exp\_13 ordering}) resolved the error onto hierarchy levels and watched it move. It
found nothing --- and, per the methodological correction above, \emph{could not have} answered
the question. It is retained only as a record of that.

\textbf{The correct probe} (\code{exp\_15 alignment}) trains a denoiser on unfiltered data and
tracks which mismatched oracle $\mathrm{BP}_k$ it is closest to. Result: \textbf{a weak
version of the effect, not the staircase.} The argmin moves in the predicted direction at the
two intermediate noise levels --- $2\to2\to0$ at $t=0.4$ and $2\to1\to1$ at $t=0.8$ over the
first ${\sim}1000$ gradient steps --- then \textbf{saturates for the remaining 31{,}000}. At
$t=0.1$, $0.2$ and $1.6$ it never leaves $k=0$.

So the direction is right, the effect is early and small, and nothing like the clean staircase
P1 reports. Three plausible reasons, none verified here: P1's oracles are far more
distinguishable (a discrete non-overlapping transition tensor versus a Gaussian edge); a
denoising regression loss converges in far fewer steps than their masked-language objective,
which needs ${\sim}10^3$ epochs; and the architecture and task differ.
\textbf{The honest summary is that this setting is too easy to resolve the phenomenon}, not
that the phenomenon is absent.

\medskip
\textbf{The per-level error is flat} across levels \emph{and} across training budget
(250 $\to$ 16{,}000 gradient steps), consistent with the above.

\textbf{But the level-resolved error is decisive about something else.} Against the right
null --- a method with uniform per-mode error --- the concentration ratio spreads across
levels by \textbf{1.6--2.2$\times$ for the $\epsilon$-network} (flat: as if the hierarchy did
not exist), 2.8--12.6$\times$ for the $x_0$-network, and \textbf{224--6957$\times$ for
EM-BP}, migrating from the finest levels at small $t$ to the coarsest mode at large $t$.
That is a quantitative sense in which BP inherits the hierarchy from the graph and the
network does not.

\emph{Methodological note already paid for:} per-level error must be reported
\textbf{absolutely}, not only relatively. A relative error divides by the reference magnitude
in that subspace, which shrinks with the level's eigenvalue, so fine levels look worse for
\emph{every} method --- including one with uniform absolute error. The first version of this
measurement showed a clean coarse-to-fine gradient that was entirely the normalisation.

\subsubsection*{Filtering removes rungs from the ladder --- a derivation worth keeping}

This is where P1's knob and P2's time scale meet in one formula, and it is the most
self-contained new piece of mathematics in Layer 6. Under the \emph{correct} filtering, the
leaf covariance is block-diagonal-plus-constant: the depth-$(L-k)$ subtree covariance inside
each of the $b^k$ blocks, and $\rho^{2L}$ between them. Its eigenvectors are

\begin{enumerate}
\item \textbf{within-block contrasts} --- the subtree's own non-uniform modes, each
multiplicity multiplied by $b^k$;
\item \textbf{between-block contrasts}, constant on each block and summing to zero across
blocks, with the single eigenvalue
$\Lambda_{\rm between} = S_0^{(L-k)} - b^{L-k}\rho^{2L}$, multiplicity $b^k - 1$;
\item the \textbf{uniform mode}, $\Lambda = S_0^{(L-k)} + (b^k-1)b^{L-k}\rho^{2L}$.
\end{enumerate}

\noindent
So \textbf{the top $k$ rungs merge into one}, and the number of distinct speciation times
falls from $L+1$ to $L-k+2$, reaching 2 at $k=L$ where the covariance is equicorrelated.
Measured at $L=4$: \textbf{5, 5, 4, 3, 2} rungs for $k=0\dots4$.

Two checks fall out for free. The analytic spectrum matches \code{eigh} of the dense
covariance to $\mathbf{10^{-15}}$ at every $(L,k)$ tried, and $\mathrm{BP}_k$ matches a dense
linear solve to $\mathbf{10^{-15}}$. And $k=0$ and $k=1$ giving the \emph{same} count
independently reproduces P1's remark that those two cases share a tree topology and differ
only in the top transition probabilities --- which in the Gaussian case means they do not
differ at all.

\subsubsection*{P6. Does the advantage grow with correlation range? --- \partialx: confounded}

\begin{center}
\begin{tabular}{@{}rrrrrr@{}}
\toprule
$k$ & block & $\hat\rho$ & EM-BP abs err & network abs err & advantage \\
\midrule
1 & 2  & 0.8796 & 0.0154 & 0.4143 & 28.7$\times$ \\
2 & 4  & 0.8926 & 0.0085 & 0.3702 & 49.2$\times$ \\
3 & 8  & 0.9010 & 0.0038 & 0.3640 & 99.0$\times$ \\
4 & 16 & 0.9041 & 0.0084 & 0.3470 & 39.4$\times$ \\
\bottomrule
\end{tabular}
\end{center}

\noindent
\textbf{The confound, which is why this is partial.} Fixing the number of \emph{sequences}
does not fix the number of \emph{edges}: a $k$-filtered depth-4 tree has $32 - 32/2^k$ edges,
i.e.\ 16, 24, 28, 30. Larger $k$ therefore hands EM more data at the same nominal budget.
The sweep moves two things at once, so the \emph{rising} advantage cannot be attributed to
correlation range. What the table supports without qualification is the negative: the
advantage \emph{does not shrink} as correlations reach further.

\subsection{[N] New questions opened by the answers}

\subsubsection*{N1. What scalar governs locality? --- \done: negentropy, not kurtosis}

\begin{sloppypar}
Candidates screened against the measured locality excess: $|$excess kurtosis$|$, negentropy
($\mathrm{KL}(f\|N(0,q))$), $L^1$ and $L^\infty$ distance to that Gaussian, and mode count.
Negentropy wins decisively --- pooled Pearson $\mathbf{+0.738}$ (Spearman $+0.842$) against
$-0.043$ for $|$kurtosis$|$.
\end{sloppypar}

\begin{center}
\begin{tabular}{@{}lrrr@{}}
\toprule
family & negentropy & $|$excess kurt$|$ & locality excess \\
\midrule
gauss\_mix $\kappa=0.6$ & 0.026 & 0.72 & 1.12 \\
Student-$t$ $\nu=5$ & 0.045 & \textbf{6.00} & 1.25 \\
laplace & 0.072 & 3.00 & 1.62 \\
uniform & 0.177 & 1.20 & 1.46 \\
gauss\_mix $\kappa=0.9$ & \textbf{0.462} & 1.62 & \textbf{4.94} \\
\bottomrule
\end{tabular}
\end{center}

\noindent
Student-$t$ has the largest kurtosis of any family and the second-\emph{smallest} negentropy:
heavy tails carry huge fourth moments but very little probability mass, so they sit close to
a Gaussian in KL. \textbf{Bimodality moves mass}, so $\kappa=0.9$ is 10$\times$ further in KL
than Student-$t$ while having a quarter of its kurtosis. That locality --- an information-flow
property --- is governed by an information-theoretic distance rather than a moment is the
natural answer in hindsight, and it is a \emph{different} scalar from the one governing
closure error, where kurtosis works well enough to extrapolate (R3). Two mechanisms, two
summaries.

\subsubsection*{N2. Does a structured architecture close the gap? --- \done: most of it, but the rest grows}
\label{sec:receptive}

This was the most important \emph{unasked} question in the project --- the one a reviewer
asks first. 4 seeds, oracle over both receptive-field radius and parameterisation for both
networks:

\begin{center}
\begin{tabular}{@{}rrrr@{}}
\toprule
$N$ & \textbf{EM-BP} (13 params) & global MLP (25{,}505) & local CNN, oracle $r$ (6{,}081) \\
\midrule
128  & $\mathbf{0.0368 \pm 0.0067}$ & $0.3595 \pm 0.0044$ & $0.0774 \pm 0.0024$ \\
512  & $\mathbf{0.0174 \pm 0.0017}$ & $0.1794 \pm 0.0035$ & $0.0570 \pm 0.0007$ \\
2048 & $\mathbf{0.0123 \pm 0.0009}$ & $0.1246 \pm 0.0027$ & $0.0505 \pm 0.0015$ \\
\bottomrule
\end{tabular}
\end{center}

\begin{center}
\begin{tabular}{@{}rrr@{}}
\toprule
$N$ & vs global MLP & vs local CNN \\
\midrule
128  & $9.76 \pm 1.79$ & $2.10 \pm 0.39$ \\
512  & $10.29 \pm 1.02$ & $3.27 \pm 0.32$ \\
2048 & $10.15 \pm 0.75$ & $\mathbf{4.11 \pm 0.32}$ \\
\bottomrule
\end{tabular}
\end{center}

\noindent
Two things this settles, one in each direction.

\textbf{The structured architecture really does close most of the gap.} Against the vanilla
MLP the margin is ${\sim}10\times$ and flat in $N$; against the locality-respecting CNN it is
2--4$\times$. Roughly \textbf{60\% of the vanilla deficit was the architecture}, not the
estimator.

\textbf{But the remaining margin \emph{grows} with data rather than closing.}
$2.10 \to 3.27 \to 4.11$, well outside the standard errors. The mechanism is visible in the
columns: EM-BP falls $0.0368 \to 0.0123$ (a factor 3.0 over 16$\times$ data) while the CNN
falls $0.0774 \to 0.0505$ (a factor 1.5) and is visibly flattening. The structured network
converges to a floor set by its \emph{approximation} error; the estimator keeps paying down
\emph{parametric} error.

\begin{keybox}
``More data will close the gap'' is the natural guess and \textbf{it is wrong here}. The
honest headline against a well-chosen architecture is \textbf{4.1$\times$ at $N=2048$ and
rising} --- not the $10\times$ of the vanilla comparison.
\end{keybox}

\subsubsection*{The optimal receptive field is finite, and exceeding it hurts}

A separate finding from the same experiment, worth stating because it corrects the project's
own earlier reading. Ledger item B15 read the Gaussian locality law as ``a local head is
near-optimal once $r \gtrsim \xi\log(1/\epsilon)$'' --- a \emph{lower} bound on the radius.
The sweep says something sharper: the error is \textbf{U-shaped in $r$}.

\begin{center}
\begin{tabular}{@{}lrrrrrrr@{}}
\toprule
$r$ & 1 & 2 & \textbf{3} & 4 & 8 & 12 & 16 \\
\midrule
error & 0.229 & 0.100 & \textbf{0.065} & 0.087 & 0.094 & 0.125 & 0.125 \\
\bottomrule
\end{tabular}
\end{center}

\noindent
(Gaussian chain, $t=0.1$, interior sites.) Doubling the radius past the optimum roughly
doubles the error --- an ordinary bias-variance trade. The optimum \emph{grows with noise
level} (3, 3, 6, 12, 16 for $t=0.1 \dots 1.6$), which the locality law predicts, and
non-Gaussian families need a \emph{wider} field than the Gaussian at matched $(\rho,t)$,
in the direction and roughly the magnitude F5 predicted.

%======================================================================
\newpage
\section{The experiment catalogue}
\label{sec:experiments}

All under \code{research/nongaussian-bp/} on the research branch. Each has a self-contained
docstring; each writes to \code{outputs/<name>/}.

\begin{center}
\small
\begin{longtable}{@{}p{2.9cm}p{9.0cm}p{2.9cm}@{}}
\toprule
\textbf{Experiment} & \textbf{What it establishes} & \textbf{Answers} \\
\midrule
\endhead
\code{exp\_01\_grid\_} \code{validation} & Grid BP as calibrated numerical ground truth against the known Gaussian score. Worst rel.\ error $9.2\times10^{-15}$ at $M=401$, $A=8$. & R1, C-iii, B10 \\
\code{exp\_02\_laplace\_} \code{gaussian\_message} & Gaussian-message error on the Laplace chain, done right. Median rel.\ score error $0.39$ at $t=0.02 \to 7.6\times10^{-5}$ at $t=2.4$. & C-i, C-iii \\
\code{exp\_03\_nongaussian\_} \code{innovation\_sweep} & Six innovation families at identical second-order structure. Noise level dominates; error \emph{decreases} with $\rho$. & R2, R3 \\
\code{exp\_04\_approx\_} \code{markovianity} & Chain $+$ global latent. The non-Markov correction is \emph{exactly rank one}. Residual MLP on frozen BP ${\sim}20\times$ more efficient. & F3, R5 \\
\code{exp\_05\_reverse\_} \code{dynamics} & Reverse SDE under exact / closed / truncated scores. Gaussian score washes out heavy tails: kurtosis 0.12 vs true 2.7--2.9. & \textbf{C-viii}, F4 \\
\code{exp\_06\_em\_parameter\_} \code{recovery} & EM correctness: monotonicity, Fisher information, misspecification, rates, quantization. & E1, L2, L3, R1 \\
\code{exp\_07\_em\_vs\_} \code{score\_network} & \textbf{The headline.} EM-BP vs a score network: sample efficiency, capacity control, transfer across $t$, inference cost. & \textbf{E3}, C-ix \\
\code{exp\_08\_gradient\_vs\_} \code{exact\_mstep} & The literal gradient route vs the exact M-step. Split is not the intuitive one. & \textbf{L1} \\
\code{exp\_09\_mixture\_} \code{message\_closure} & Mixture \emph{message} closure. Separates representation error from model error for the first time. & \textbf{R4/F1}, R3 \\
\code{exp\_10\_discrete\_} \code{alphabet} & Discrete alphabet: headline retested with \emph{zero} grid error. Baum-Welch identity. & \textbf{F2} \\
\code{exp\_11\_nongaussian\_} \code{locality} & Is the locality law second-order universal? No --- rate is not, excess up to 2.45$\times$. & \textbf{F5}, N1 \\
\code{exp\_12\_receptive\_} \code{field} & How big a receptive field a score head needs. U-shaped in $r$; CNN closes 60\% of the gap. & \textbf{N2}, C-x, R5 \\
\code{exp\_13\_speciation\_} \code{cascade} & Hierarchical priors: the speciation ladder, and where each method's error sits. & \textbf{P1}, P2, P5 \\
\code{exp\_14\_memorization\_} \code{collapse} & Collapse/memorisation, and why a BP score cannot collapse. & \textbf{P3} \\
\code{exp\_15\_bp\_oracles} & Which BP oracle a learned diffusion score agrees with. & P5 \\
\bottomrule
\end{longtable}
\end{center}

\subsection{Tests}

101 tests across 10 files in \code{research/nongaussian-bp/tests/}, all passing:

\begin{center}
\small
\begin{tabular}{@{}ll@{}}
\toprule
\code{test\_em\_bp.py} (23 tests) & \code{test\_hierarchy.py} (51 tests) \\
\code{test\_bp\_mixture.py} & \code{test\_discrete.py} \\
\code{test\_gaussian\_bp\_equivalence.py} & \code{test\_gaussian\_precision\_score.py} \\
\code{test\_grid\_convergence.py} & \code{test\_local\_head.py} \\
\code{test\_message\_normalization.py} & \code{test\_score\_mean\_error\_identity.py} \\
\bottomrule
\end{tabular}
\end{center}

\noindent
Separately, the earlier analytic packages carry their own audits, both re-executed
2026-07-14:

\begin{center}
\small
\begin{tabular}{@{}ll@{}}
\toprule
\code{research/gaussian-bp/code/numerical\_audit.py} & \textbf{72/72 pass} \\
\code{research/unified-note/code/numerical\_audit.py} & \textbf{55/55 pass} \\
\bottomrule
\end{tabular}
\end{center}

%======================================================================
\newpage
\section{The cluster sweep}
\label{sec:cluster}

Run on the Bocconi HPC, 30 July -- 4 August 2026. Three sweeps, \textbf{5{,}530 rows}, zero
failed cells, 48 CPU-hours. This is the \emph{separate} implementation
(\S\ref{sec:duplication}), so it functions as independent replication.

\begin{center}
\begin{tabular}{@{}lrl@{}}
\toprule
sweep & rows & design \\
\midrule
\code{identifiability} & 4{,}320 & 5 families $\times$ 4 $\rho$ $\times$ 3 $n$ $\times$ 8 $t$ $\times$ 3 $M$ $\times$ 3 reps \\
\code{sample\_efficiency} & 1{,}050 & 5 families $\times$ 2 $\rho$ $\times$ 7 $M$ $\times$ 3 $t_{\rm eval}$ $\times$ 5 reps \\
\code{nonlinearity\_gap} & 160 & 5 families $\times$ 4 $\rho$ $\times$ 8 $t$ \\
\bottomrule
\end{tabular}
\end{center}

\subsection{The closure gap is largest at \emph{small} $t$}

Relative score error of the Gaussian closure against exact grid BP, $\rho=0.85$, $n=40$:

\begin{center}
\begin{tabular}{@{}lrrrrr@{}}
\toprule
$t$ & gaussian & uniform & laplace & mixture & student \\
\midrule
0.08 & $7.0\times10^{-16}$ & 0.176 & \textbf{0.200} & 0.140 & 0.180 \\
0.15 & $4.0\times10^{-16}$ & 0.092 & 0.127 & 0.069 & 0.108 \\
0.30 & $2.7\times10^{-16}$ & 0.036 & 0.060 & 0.027 & 0.056 \\
0.50 & $1.8\times10^{-16}$ & 0.015 & 0.030 & 0.013 & 0.032 \\
0.80 & $1.2\times10^{-16}$ & 0.006 & 0.013 & 0.005 & 0.013 \\
1.20 & $8.4\times10^{-17}$ & 0.002 & 0.005 & 0.002 & 0.006 \\
1.80 & $5.5\times10^{-17}$ & $3.5\times10^{-4}$ & $7.9\times10^{-4}$ & $2.6\times10^{-4}$ & $1.2\times10^{-3}$ \\
2.40 & $4.4\times10^{-17}$ & $3.8\times10^{-5}$ & $9.2\times10^{-5}$ & $2.7\times10^{-5}$ & $1.5\times10^{-4}$ \\
\bottomrule
\end{tabular}
\end{center}

\noindent
The \textbf{gaussian column is Proposition~C-v holding at machine precision} across 32
independent cells. The gap decays \emph{monotonically} in $t$ for every family --- confirming
on separate code that the rise-with-$t$ J\'er\^ome found suspicious in the call was the bug,
not physics. Peak gap over the whole sweep: 0.46 (uniform, $\rho=0.5$, $t=0.08$).

\begin{keybox}
\textbf{Interpretation, and why it makes C-viii \emph{more} urgent.} Large $t$ Gaussianises
everything, so the closure becomes exact; the prior's non-Gaussianity survives only where
the likelihood is weak, i.e.\ small $t$. But J\'er\^ome's own argument in the call is that
small-$t$ score errors are \emph{cheap} --- $\alpha_t\approx 1$, the denoiser barely acts.
So we do not currently know whether the 20\% closure gap costs anything in generated
samples. Only reverse diffusion under a learned score can say.
\end{keybox}

\subsection{EM parameter recovery, and where it breaks}

Median $|\hat\rho - \rho|$, pooled over families, $\rho$, $n$, reps:

\begin{center}
\begin{tabular}{@{}lrrrr@{}}
\toprule
$t$ & $M=250$ & $M=1000$ & $M=4000$ & EM converged \\
\midrule
0.08 & 0.0039 & 0.0021 & 0.0009 & 100\% \\
0.15 & 0.0046 & 0.0024 & 0.0011 & 100\% \\
0.30 & 0.0059 & 0.0028 & 0.0013 & 100\% \\
0.50 & 0.0071 & 0.0036 & 0.0020 & 95\% \\
0.80 & 0.0102 & 0.0075 & 0.0037 & 52\% \\
1.20 & 0.0284 & 0.0243 & 0.0211 & 16\% \\
1.80 & 0.0961 & 0.0749 & 0.0725 & 11\% \\
2.40 & 0.1563 & 0.1422 & 0.1232 & 27\% \\
\bottomrule
\end{tabular}
\end{center}

\noindent
More data barely helps past $t=1.2$. That signature --- error plateaus in $M$, convergence
flag collapses --- is an \textbf{optimiser} failure, not an information failure, and it is
confirmed directly: at $M=100{,}000$, $t=1.8$, EM from a random start lands \textbf{735.6
nats worse in total} than EM from the truth, and only 2 of 24 restarts find the right basin.
A moment-based initialiser fixes this up to $t\approx 2.4$ and then fails honestly at
$t=3.5$, where the signal $\alpha^2 = 9.1\times10^{-4}$ really is below the sampling error
$3.5\times10^{-3}$. \emph{Two different regimes, both real.}

\begin{keybox}
\textbf{Not usable from this sweep:} \code{ridge\_width} is censored in 1{,}889 of 4{,}320
cells (44\%). Do not quote it. The censoring flag is in the CSV.
\end{keybox}

%======================================================================
\newpage
\section{Errors found and corrected}
\label{sec:corrections}

This section exists because a research record that lists only successes cannot be trusted.
Each of these was found \emph{after} a result had been reported from it.

\begin{center}
\small
\begin{longtable}{@{}p{4.3cm}p{6.4cm}p{4.2cm}@{}}
\toprule
\textbf{What was wrong} & \textbf{Consequence} & \textbf{Resolution} \\
\midrule
\endhead
\code{rng\_for} not reproducible across processes & Seeded from Python's builtin \code{hash}, salted per process (PEP 456). \textbf{No experiment in the package was bit-reproducible from a fresh interpreter}, contrary to the README. & blake2b digest. Pairing \emph{within} a run held, so exp\_01--05 numbers stand as measurements. \\
\addlinespace
Score-network baseline handicapped & It predicted the clean signal only, whereas the standard recipe predicts the \emph{noise}. Not equivalent in finite samples. The headline rested on a crippled competitor. & Both parameterisations now trained and reported; $\epsilon$ wins at small $t$, $x_0$ at large $t$, so the result rests on neither. \\
\addlinespace
MDN M-step broke EM's guarantee & A fixed Adam step is not an ascent step: likelihood rose overall while violating monotonicity by 547 nats. & Backtrack until $Q$ actually increases. Zero violation, negligible cost. \\
\addlinespace
exp\_06 Part 5 seeding confound & Seed keyed on grid size, so every cell drew different data. Reported ``at $M=201$ three true values all return 0.750''. & Re-run with data held fixed: the collapse is real but lands on \textbf{0.800}. The phenomenon survived; that number did not. \\
\addlinespace
$N^{-1/2}$ rate declared unsupported & 4 replicates gave slopes $-0.26,-0.69,-0.25,-0.41$; conclusion ``rate not established''. Correct \emph{for that data}. & 12 replicates $\Rightarrow$ $-0.500\pm0.048$. The lesson is about replicate counts. \\
\addlinespace
Per-level error normalisation & Relative error divides by a magnitude that shrinks with level, so fine levels look worse for \emph{every} method. First version showed a clean coarse-to-fine gradient that was \textbf{entirely the normalisation}. & All three forms (relative, absolute, share of squared error) now written. \\
\addlinespace
exp\_13 \code{ordering} probe & Wrong measurement for the question. & Replaced by \code{exp\_15}, which uses the probe the paper actually uses. \\
\addlinespace
Transfer-across-$t$ premise & The experiment was designed around ``the network must extrapolate in $t$''. There is \textbf{no cliff} at the schedule boundary. & The extrapolation claim was withdrawn. The difficulty is in $x$-space, not $t$-space. \\
\addlinespace
Single-seed architecture figures & 3.3$\times$ / 2.75$\times$ quoted from one seed at $N=1024$. & Replaced by the 4-seed oracle-fair run: 4.11$\times$ at $N=2048$. \\
\addlinespace
P1's filtering construction & Implemented as $b^k$ \emph{independent} subtrees. The paper draws depth-$k$ nodes conditionally independently \emph{given the root}, so blocks stay correlated through it. Independent blocks put a \textbf{zero} where the truth puts $\rho^{2L}$. & Found by reading P1 in full. Reimplemented; a test now pins the distinction. \\
\addlinespace
``The hierarchy $\times$ diffusion intersection is untouched'' & False. Sclocchi, Favero and Wyart (PNAS 2025, arXiv:2402.16991), \textbf{cited by P1 itself}, studies exactly that. & Claim withdrawn. The measurements are \emph{consistent with} that work, not novel against it --- which is more useful than what it replaced. \\
\bottomrule
\end{longtable}
\end{center}

\subsection{Two artifacts that are real, not bugs}

\textbf{The Laplace M-step is exact but lattice-quantized} (L3). $4/5$ is an attractor;
$\rho^*=0.7913$ pins at $0.8000$ across an 8$\times$ grid refinement with constant $0.0087$
bias. Real for the discretized model, a \emph{bias} rather than a resolution limit.
Consequence: every rate claim uses the Gaussian or mixture kernel.

\textbf{More expressive is not better.} At $N=400$ the neural kernel (873 params) is dominated
on \emph{both} axes by the 13-parameter mixture kernel: training log-likelihood $-16781$ vs
$-16656$, held-out score error $0.052$ vs $0.026$. The efficiency argument applies to the
structured method's own internals, not only to the baseline it is compared against.

%======================================================================
\newpage
\section{Limitations}
\label{sec:limitations}

Stated as they are recorded in the package, not softened.

\begin{enumerate}
\item \textbf{Inference cost is the honest loss.} Grid BP is $O(nM^2)$ per chain per
evaluation; a network forward pass is a few matmuls. Measured slowdown \textbf{211--320$\times$}
at $M=401$. Since reverse diffusion evaluates the denoiser at every integration step, this is
not a detail. Distilling the fitted BP denoiser into a network, or using Gaussian-projected
BP at inference, is the obvious remedy --- \emph{not something claimed here}.

\item \textbf{EM is slower to train here.} 443\,s at $N=2048$ vs ${\sim}37$\,s for the
network, and EM's cost grows linearly in $N$ while the network's is fixed by its step count.
\textbf{The win is in data, not wall-clock.}

\item \textbf{Everything rests on the chain assumption.} If the prior is not Markov the
estimator is misspecified in a way no amount of data fixes. Layer 4 showed the correction is
exactly rank one for AR(1)$+$global-latent priors, suggesting a hybrid --- untested (F3).

\item \textbf{The margin against a well-chosen architecture is 4$\times$, not 10$\times$.}
The 10$\times$ figure is against a vanilla MLP as the email specified. See N2.

\item \textbf{One chain family, $n=32$.} Linear AR(1) with non-Gaussian innovations. The
claim is demonstrated in this setting, not in general.

\item \textbf{The chain has no speciation cascade} (P2). Its top eigenvalue is bounded, so
results here do not transfer to hierarchically structured data without an explicit argument.

\item \textbf{The Layer-6 paper connections were established without reading the papers}
(access blocked). Everything is re-derived and verified internally, but the correspondence to
what the papers actually claim is unverified.

\item \textbf{The sample-efficiency advantage is tied to small parametric dimension} (F2). At
$S=12$ / 132 parameters it is down to 4.4$\times$ and falling.
\end{enumerate}

%======================================================================
\newpage
\section{What is missing, in priority order}
\label{sec:missing}

\begin{center}
\begin{tabular}{@{}p{0.6cm}p{5.4cm}p{6.4cm}p{2.6cm}@{}}
\toprule
& \textbf{Item} & \textbf{Why it matters} & \textbf{Cost} \\
\midrule
1 & \textbf{The paper / Overleaf} (E4, C-xii) & The only item with the 16 Sep deadline. The
theory document exists and is paper-shaped; it needs a results section and the narrative arc
J\'er\^ome supplied. & days, no compute \\
\addlinespace
2 & \textbf{Reverse diffusion under a \emph{learned} score} (C-viii, F4) & The one metric that
says whether \emph{any} of the score gaps matter. Unblocked now that Layer 5 produces a
learned score. Also where the 211--320$\times$ inference cost bites. & ${\sim}1$ day $+$ compute \\
\addlinespace
3 & \textbf{Learned kernel $+$ rank-one correction} (F3) & Turns ``the prior is Markov'' from
an assumption into an approximation with a measured correction. The natural product of Layers
4 and 5. & compute \\
\addlinespace
4 & \textbf{The literal depth question} (C-x) & J\'er\^ome's stated ``biggest question'': can
a narrow, deep network implement the two BP sweeps, and in how many layers? exp\_12 answered
receptive field, not depth. & open-ended \\
\addlinespace
5 & Distillation of the BP denoiser & Directly addresses limitation 1. & ${\sim}1$ day \\
\addlinespace
6 & Nonparametric collapse test & Answers the obvious objection to the P3 memorisation result. & compute \\
\addlinespace
7 & P6 with edge count fixed & Removes the confound; would let the correlation-range claim be
made. & compute \\
\addlinespace
8 & \textbf{Read P2 and Sclocchi--Favero--Wyart} & P1 is read; P2 is not, and a write-up
citing it must use its symbols. Sclocchi et al.\ (arXiv:2402.16991) is the closest prior work
on hierarchy $\times$ diffusion. & hours \\
\bottomrule
\end{tabular}
\end{center}

\subsection{The Layer-6 package's own ranked list}

Reproduced from \code{docs/PAPER\_CONNECTIONS.md} §6, because it is more specific than the
table above about what would move Layer 6 forward:

\begin{enumerate}
\item \textbf{Read P2}, and \textbf{read Sclocchi, Favero and Wyart (arXiv:2402.16991)}
before writing anything about hierarchy and diffusion.
\item \textbf{The nonparametric collapse test.} The obvious objection to P3 is that a
two-parameter kernel \emph{cannot} memorise. \code{exp\_14 --set include\_mdn=True} runs the
same comparison with a mixture-density kernel of several hundred parameters and the same
$n$-independent statistic, which \textbf{separates the mechanism from the parameter count}.
\item \textbf{A discrete-alphabet tree} --- closest to the actual data model of the
hierarchical-filtering paper. \code{src/discrete.py} does chains, \code{src/hierarchy.py}
does continuous trees; \textbf{the two have never been crossed}.
\item \textbf{The cascade under a learned score} rather than the exact one: which levels of
the ladder each method gets right \emph{dynamically}.
\item \textbf{A non-Gaussian tree.} \code{tree\_bp\_grid} and \code{fit\_em\_tree} already
support it with no new machinery; the Layer-3 question --- recover an unknown innovation law
from noisy data alone --- transfers directly.
\end{enumerate}

\noindent
Item 2 is the one a referee asks first, and it is a one-flag run.

\subsection{Repository housekeeping}

\begin{itemize}
\item \textbf{Merge the branch, or at least stop working off \code{main}.} 49 commits of
research sit unmerged. This is how the thread was lost.
\item \code{research/em-learning/} is untracked and duplicative. Keep \code{outputs/} as a
replication artifact; the code is superseded.
\item Stale docstrings in \code{research/em-learning/code/}: \code{bp\_core.py} cites a
non-existent \code{analytics.py}; \code{chain\_models.py} cites a non-existent
\code{sampling.py}; \code{run\_sweep.py} claims a network ``cannot match without
retraining'', which is false and which \code{nn\_baseline.py}'s own docstring contradicts.
\end{itemize}

%======================================================================
\newpage
\section{How to reproduce anything here}

\subsection{Get the code}

\begin{keybox}
\code{cd /Users/gloriabagnato/Code/Thesis/Diffusion}\\
\code{git fetch origin}\\
\code{git checkout claude/em-bp-denoiser-learning-e07ike}
\end{keybox}

\subsection{Run an experiment}

Each experiment is standalone and writes to \code{outputs/<name>/}:

\begin{keybox}
\code{cd research/nongaussian-bp}\\
\code{python experiments/exp\_07\_em\_vs\_score\_network.py}
\end{keybox}

\subsection{Run the tests}

\begin{keybox}
\code{cd research/nongaussian-bp \&\& python -m pytest tests/ -q}
\end{keybox}

\noindent
Expect 101 passing. The five cross-checks of L2 are permanent tests, not one-off scripts.

\subsection{The cluster}

\begin{keybox}
\code{ssh lnode01-da.hpc.unibocconi.it}\\
\code{cd \textasciitilde/bp-em-diffusion}\\
\code{sbatch --array=0-19 code/sweep.sbatch}\quad\emph{(QOS cap: 30 submitted, 10 running)}
\end{keybox}

\noindent
Validation runs first with \code{--dependency=afterok}, so bad foundations cannot burn
compute. Sharding shuffles the cell list with a fixed seed (12345) before the modulo split,
because the raw ordering has period-9 inner loops that alias badly against shard counts.

\subsection{Build the reports}

\begin{keybox}
\code{cd research/nongaussian-bp/report \&\& latexmk -pdf em\_bp\_learning.tex}
\end{keybox}

\noindent
(Or \code{tectonic em\_bp\_learning.tex} --- no TeX Live is installed on the laptop.)

\vspace{2em}
\begin{keybox}
\textbf{Provenance of this document.} Compiled 4 August 2026 from:
\code{questions/QUESTIONS\_AND\_ANSWERS.md} (798 lines),
\code{docs/EM\_BP\_LEARNING\_COMPENDIUM.md} (872 lines),
\code{docs/PAPER\_CONNECTIONS.md} (499 lines),
\code{docs/RESULT\_LEDGER.md}, \code{docs/AGENT\_HANDOFF\_EM\_BP.md},
\code{report/layer6\_session\_report.tex}, the experiment docstrings,
\code{research/nongaussian-bp/hpc/README.md}, and the merged cluster CSVs.

\medskip
\textbf{Scope.} This covers the 49 commits that constitute the
\code{claude/em-bp-denoiser-learning-e07ike} branch --- 175 files, 22{,}749 insertions,
essentially all of \code{research/nongaussian-bp/} plus the four \code{docs/} files. Work
that predates the branch and already sits on \code{main} --- the thesis chapters, the
Gaussian AR(1) analytic packages, the AMP/TAP appendix, the conjecture-15 experiments --- is
referenced where relevant but is \emph{not} catalogued here.
The transcript of the 29 July call and the PDF of Marc's email were re-read in full for the
question inventory. Numbers are quoted from those sources; where this document states a
limitation or a correction, it is because the source states it, not because it was inferred.
\end{keybox}

\end{document}
